\chapter{Revenue and Profit across the Industry}\label{ch:in:revenue-and-profit-across-the-industry}

Put every firm that files side by side and revenue per employee ranges over more than a factor of twenty. In 2025
dollars, one exchange group took in \$0.48 million a head in 2023, and one London systematic investment manager \$13.0
million a head in its year to January 2022. Within a firm it moves almost as much: a listed market maker's \omterm{def:in:revenue-and-profit-across-the-industry:perhead}{revenue per
head} was 4.5 times as high in 2020, the most volatile year of the period, as in 2023, one of the quietest. The table
this chapter builds is the industry's income as the filings show it, and its gaps are part of the result: the firms
that file nothing include most of the private market makers that the earlier chapters described. The chapter builds
the table from sources, says what a per-head figure measures and hides, measures how \omterm{def:in:revenue-and-profit-across-the-industry:perhead}{revenue per head} moves with the
market's volatility for each kind of firm, and names who is missing.

\section{A sourced table of filed numbers}

\begin{definition}[Revenue per head, profit per head]\label{def:in:revenue-and-profit-across-the-industry:perhead}
\emph{Revenue per head}\index{revenue per head} is a firm's revenue for a year, after the costs that pass straight
through it to others, divided by its headcount; \emph{profit per head}\index{profit per head} is a stated profit line
(operating profit, profit before tax, or net profit) divided by the same headcount. Both are reported with the revenue
or profit line, the basis of the headcount (at the year end, or the year's average), the currency and the price level
used.
\end{definition}

Each part of the definition is a choice that changes the answer, so the table records all of them. The revenue line is
the filer's top line after pass-through costs, because a gross line inflates the figure for firms that collect money
for others. An exchange group that reports total revenues of \$12\,640 million for 2025 also reports
transaction-based expenses of \$2\,709 million: \$412 million of regulatory (section~31) fees it collects and passes on,
and \$2\,297 million of cash liquidity payments, routing and clearing. Its revenue after those, \$9\,931 million, equals its operating expenses plus its operating income, which is how the table
derives it from the tagged facts. A market maker's revenue is taken after the brokerage, exchange and clearance fees and
the financing costs that Book~16, chapter~1, calls volume-driven; a bank's and a broker's after interest expense (net
revenues); a fund manager's is its fees. The headcount is whatever the filing states: most US filers give a count at
the year end, UK companies the year's monthly average, one bank a figure rounded to thousands, and one European market
maker full-time equivalents.

\begin{dated}{2025-12}{Revenue per head from filings, latest year}\label{dat:in:revenue-and-profit-across-the-industry:table}
\footnotesize
\begin{tabular}{@{}llrrrr@{}}
\toprule
filer & kind & year & revenue, m & staff & \$m a head\\
\midrule
Nasdaq & exchange & 2025 & \$5\,249.0 & 9\,525 & 0.55\\
Man Group & asset manager & 2025 & \$1\,325.0 & 1\,719 & 0.77\\
Intercontinental Exchange & exchange & 2025 & \$9\,931.0 & 12\,844 & 0.77\\
Morgan Stanley & bank (firm) & 2025 & \$70\,645.0 & 83\,000 & 0.85\\
Flow Traders & market maker & 2025 & \euro485.8 & 635 & 0.86\\
MarketAxess & venue & 2024 & \$817.1 & 891 & 0.94\\
Goldman Sachs & bank (firm) & 2025 & \$58\,283.0 & 47\,400 & 1.23\\
Tradeweb & venue & 2025 & \$2\,052.4 & 1\,569 & 1.31\\
Coinbase & crypto exchange & 2025 & \$7\,181.3 & 4\,951 & 1.45\\
Cboe Global Markets & exchange & 2025 & \$2\,429.1 & 1\,661 & 1.46\\
CME Group & exchange & 2025 & \$6\,520.6 & 3\,875 & 1.68\\
Interactive Brokers & broker & 2024 & \$5\,185.0 & 2\,998 & 1.77\\
Virtu Financial & market maker & 2025 & \$2\,214.9 & 1\,027 & 2.16\\
Optiver & market maker & 2025 & \euro4\,556.0 & $>2\,000$ & $<2.57$\\
Quadrature Capital & asset manager & FY25 & \pounds1\,224.8 & 173 & 9.29\\
\bottomrule
\end{tabular}

\smallskip\noindent Revenue after pass-through costs; headcount at the year end except Quadrature (average) and
Optiver (a stated lower bound); FY25 is the year to January 2025. Per head in 2025 US dollars: ECB annual average rates and the US consumer price index.
Every cell with its filing in \texttt{data/industry/filings\_perhead.csv} and \texttt{headcounts.csv}.
\end{dated}

The fifteen filers are the panel of this chapter, 2019 to 2025 where the filings allow: six exchanges and trading
venues, three market makers, two banks as whole firms, two asset managers, one broker and one crypto exchange, with 87
firm-years in all once the years without a headcount are dropped. Nine are read from SEC company facts with
\texttt{firm.filings} (chapter~11), two from Book~16's derived tables, and the rest from annual reports, results releases
and the UK register, transcribed with their pages.

\omterm{def:in:revenue-and-profit-across-the-industry:perhead}{Profit per head} is harder to pool, because the natural profit line differs by kind of firm. Four examples, each on its
own line: Quadrature's operating profit was \pounds3.22 million a head in its year to January 2025 (\$4.11 million);
CME Group's operating income \$1.09 million in 2025; Virtu's income before taxes \$1.07 million in 2025; Optiver's net
profit \euro0.65 million in 2024 (\$0.71 million), on its stated 2\,100 employees.

\section{Per head: what it measures and what it hides}

\omterm{def:in:revenue-and-profit-across-the-industry:perhead}{Revenue per head} measures how much income a firm's people produce between them, at the firm's own prices, in its own
year. It is not pay, and it is not productivity in any sense a manager could act on. What it hides is best seen firm by
firm (\cref{fig:in:revenue-and-profit-across-the-industry:ranges}).

\begin{omfigure}
\begin{tikzpicture}
\begin{axis}[width=10.5cm,height=9cm,xmode=log,xmin=0.3,xmax=20,ymin=0.4,ymax=16.6,
  ytick=data,yticklabels from table={figdata/industry/12-revenue-and-profit-across-the-industry/firm_ranges.csv}{firm},
  yticklabel style={font=\scriptsize},xticklabel style={font=\footnotesize},
  log ticks with fixed point,xtick={0.5,1,2,5,10},xlabel={\small revenue per head (\$ million of 2025)},
  grid=major,grid style={gray!20},table/col sep=comma]
\addplot+[only marks,mark=*,mark size=1.6pt,color=omDef,error bars/.cd,x dir=both,x explicit,
  error bar style={line width=1.2pt,omDef!45}] table[x=last,y=pos,
  x error plus expr=\thisrow{hi}-\thisrow{last},x error minus expr=\thisrow{last}-\thisrow{lo}]
  {figdata/industry/12-revenue-and-profit-across-the-industry/firm_ranges.csv};
\end{axis}
\end{tikzpicture}

\omcaption{\omterm{def:in:revenue-and-profit-across-the-industry:perhead}{Revenue per head} of the panel's filers, 2019--2025, in 2025 dollars on a log scale: the bar spans each
filer's lowest and highest year, the dot marks its latest. Filers are ordered by their median year. The UK entity is a
subsidiary of a larger group, shown for comparison and never pooled. Data: \texttt{data/industry/filings\_perhead.csv},
through \texttt{in\_industry.firm\_ranges}.}\label{fig:in:revenue-and-profit-across-the-industry:ranges}
\end{omfigure}

\begin{itemize}
\item \textbf{The business model.} The exchanges' figures cluster between \$0.5 million and \$1.7 million a head and
barely move; a market maker's range is wide. Flow Traders' highest year is 4.5 times its lowest, Virtu's 2.5 times;
CME Group's 1.2 times.
\item \textbf{Acquisitions.} Nasdaq's headcount rose from 6\,377 to 8\,525 in 2023, ``primarily due to our acquisition
of Adenza'', and its \omterm{def:in:revenue-and-profit-across-the-industry:perhead}{revenue per head} fell to its lowest, \$0.48 million; the Intercontinental Exchange's rose from
8\,911 to 13\,222 the same year. An acquisition adds people at once and revenue over the following years.
\item \textbf{The whole firm, not the desk.} The two banks enter as whole firms, with their wealth, lending and
investment banking staff. Their markets divisions publish revenue (chapter~7) but no headcount, so the division's
figure cannot be computed from filings.
\item \textbf{The entity, not the firm.} Hudson River Trading's UK subsidiary had 96 employees on average in 2022 and 21
in 2025, while its revenue after transaction costs rose; its \omterm{def:in:revenue-and-profit-across-the-industry:perhead}{revenue per head} went from \$1.00 million to \$7.18 million.
It also books service income from other group companies. Chapter~11's pitfalls apply: the table keeps such entities out
of every pooled number.
\item \textbf{The count.} A count at the year end, an average and a full-time-equivalent count differ most when a firm
grows or cuts fast; ``more than 2\,000'' is a lower bound, so the \omterm{def:in:revenue-and-profit-across-the-industry:perhead}{revenue per head} it gives is an upper bound.
\item \textbf{Currency and price level.} A European firm's euro revenue converted at the year's average rate moves with
the exchange rate; and 2019 dollars bought about a quarter more than 2025 dollars, which the deflator removes.
\end{itemize}

\begin{example}[Normalising one cell]\label{ex:in:revenue-and-profit-across-the-industry:cell}
Flow Traders' net trading income for 2020 was \euro933.4 million with 554 full-time equivalents at the year end:
\euro1.68 million a head. At 2020's average rate of \$1.1422 to the euro that is \$1.92 million, and at 2025's price
level, with the consumer price index at 135.83 against 109.20, \$2.39 million. The same firm's 2023 figure, \euro300.3
million over 646, becomes \$0.53 million. Each step is small; together they decide whether two firms in two currencies
and two years can be compared at all.
\end{example}

\begin{method}[Building the industry table]\label{met:in:revenue-and-profit-across-the-industry:table}
\begin{enumerate}
\item For each filer and year, take the revenue line after pass-through costs, in the filer's currency, with its
source; derive it from tagged facts where the filer does not report it as one line.
\item Take the headcount the filing states, with its basis, and a quoted phrase as evidence.
\item Convert to dollars at the year's average rate and to the base year's price level.
\item Divide; mark lower bounds and entities; never fill a missing year.
\end{enumerate}
\end{method}

\section{Volatility regimes: how results move with the market}

The market's volatility sets how much trading there is and how wide spreads are, and so much of the revenue of
anyone who earns from trading. The chapter measures it by the yearly mean of the VIX index's daily closes (Book~16's
table): 15.39 in 2019, 29.25 in 2020, 19.66 in 2021, 25.64 in 2022, 16.85 in 2023, 15.55 in 2024 and 18.93 in 2025.
\Cref{fig:in:revenue-and-profit-across-the-industry:series} shows five filers through those years.

\begin{omfigure}
\begin{tikzpicture}
\begin{axis}[width=10.5cm,height=6cm,ymode=log,xmin=2018.7,xmax=2025.3,ymin=0.4,ymax=4,
  xtick={2019,2020,2021,2022,2023,2024,2025},xticklabel style={/pgf/number format/1000 sep={},font=\footnotesize},
  log ticks with fixed point,ytick={0.5,1,2,3},yticklabel style={font=\footnotesize},
  ylabel={\small \$ million of 2025 per head},grid=major,grid style={gray!20},table/col sep=comma,
  legend style={font=\scriptsize,legend cell align=left,at={(0.5,-0.18)},anchor=north,/tikz/every even column/.append style={column sep=6pt}},
  legend columns=3]
\addplot[omDef,thick,mark=*] table[x=year,y=Virtu]{figdata/industry/12-revenue-and-profit-across-the-industry/series.csv};
\addplot[omThm,thick,mark=square*] table[x=year,y=Flow_Traders]{figdata/industry/12-revenue-and-profit-across-the-industry/series.csv};
\addplot[black!70,thick,mark=triangle*] table[x=year,y=CME_Group]{figdata/industry/12-revenue-and-profit-across-the-industry/series.csv};
\addplot[omExo,thick,mark=diamond*] table[x=year,y=Goldman_Sachs]{figdata/industry/12-revenue-and-profit-across-the-industry/series.csv};
\addplot[gray,thick,dashed,mark=o,mark options={solid}] table[x=year,y=Nasdaq]{figdata/industry/12-revenue-and-profit-across-the-industry/series.csv};
\legend{Virtu (market maker),Flow Traders (market maker),CME Group (exchange),Goldman Sachs (bank),Nasdaq (exchange)}
\end{axis}
\end{tikzpicture}

\omcaption{\omterm{def:in:revenue-and-profit-across-the-industry:perhead}{Revenue per head} through the volatility regimes of 2019--2025 (log scale). The two market makers peak in 2020,
the year of the highest mean VIX, and fall back by 2023; the exchanges and the bank move far less. Data: as
\cref{fig:in:revenue-and-profit-across-the-industry:ranges}.}\label{fig:in:revenue-and-profit-across-the-industry:series}
\end{omfigure}

The medians by kind tell the same story with more firms behind them (whole firms only, \$ million of 2025 a head):

\begin{center}\footnotesize
\begin{tabular}{@{}lrrrrrrr@{}}
\toprule
kind & 2019 & 2020 & 2021 & 2022 & 2023 & 2024 & 2025\\
\midrule
market makers & 0.90 & 2.70 & 1.62 & 1.24 & 0.96 & 1.77 & 1.51\\
exchanges and venues & 1.31 & 1.27 & 1.23 & 1.13 & 1.05 & 1.10 & 1.31\\
banks (whole firms) & 1.20 & 1.37 & 1.28 & 0.90 & 1.08 & 0.99 & 1.04\\
\bottomrule
\end{tabular}
\end{center}

\noindent The market makers' median tripled from 2019 to 2020 and fell back by 2023; the exchanges' moved within a
band of a quarter, and the banks' was highest in 2020 and 2021.
(Two or three firms stand behind each market-maker cell, so a median there is close to a single firm's figure.)

The two market makers' \omterm{def:in:revenue-and-profit-across-the-industry:perhead}{revenue per head} in 2020 was 2.16 (Virtu) and 4.51 (Flow Traders) times their 2023 figure,
while the mean VIX was 1.74 times as high. A regression puts a number on the pattern for every kind of firm at once.

\begin{method}[Elasticity of revenue per head to volatility]\label{met:in:revenue-and-profit-across-the-industry:fe}
Fit, over all firm-years of whole firms with a point headcount,
\[
\log\left(\frac{R}{N}\right)_{it}=a_i+b_{k(i)}\log \bar V_t+\varepsilon_{it},
\]
where $R/N$ is \omterm{def:in:revenue-and-profit-across-the-industry:perhead}{revenue per head} in constant dollars, $a_i$ a fixed effect for firm $i$ (its own level), $\bar V_t$ the
year's mean VIX and $b_k$ the elasticity for firms of kind $k$. The firm effects remove every difference of level
between firms, so each $b_k$ is estimated from how each firm's own figure moves across years. Report $b_k$ with its
standard error, the number of firm-years and firms, and the share of within-firm variance explained.
\end{method}

\begin{omfigure}
\begin{tikzpicture}
\begin{axis}[width=10.5cm,height=5.6cm,xmin=0.4,xmax=6.6,ymin=-2.6,ymax=2,xtick=data,
  xticklabels from table={figdata/industry/12-revenue-and-profit-across-the-industry/elasticities.csv}{kind},
  xticklabel style={font=\scriptsize,text width=1.5cm,align=center},yticklabel style={font=\footnotesize},
  ylabel={\small elasticity to mean VIX},ytick={-2,-1,0,1,2},grid=major,grid style={gray!20},table/col sep=comma,
  extra y ticks={0},extra y tick style={grid style={black!60,thick}}]
\addplot+[only marks,mark=*,mark size=2pt,color=omDef,error bars/.cd,y dir=both,y explicit,
  error bar style={thick,omDef!60}] table[x=pos,y=b,y error=err]
  {figdata/industry/12-revenue-and-profit-across-the-industry/elasticities.csv};
\end{axis}
\end{tikzpicture}

\omcaption{Elasticity of \omterm{def:in:revenue-and-profit-across-the-industry:perhead}{revenue per head} to the yearly mean VIX, by kind of firm, with bars of two standard errors: 87
firm-years of 15 whole firms, 2019--2025, with firm fixed effects. Only the market makers' elasticity is clearly
positive. Data: as \cref{fig:in:revenue-and-profit-across-the-industry:ranges}, through
\texttt{in\_industry.fit}.}\label{fig:in:revenue-and-profit-across-the-industry:elasticity}
\end{omfigure}

The market makers' elasticity is 1.19 with a standard error of 0.26: a year whose mean VIX is 10\% higher comes with
\omterm{def:in:revenue-and-profit-across-the-industry:perhead}{revenue per head} about 12\% higher, for the same firm. For the exchanges and venues it is 0.13 (0.16), for the asset
managers 0.23 (0.27) and for the banks as whole firms 0.03 (0.32): none distinguishable from zero. The broker's is
$-0.39$ (0.37), and the crypto exchange's $-1.27$ (0.59), because its revenue follows the prices of crypto assets, which
fell hardest in 2022, a year of high equity volatility (chapter~10). The slopes explain 30\% of the variance within
firms (\cref{fig:in:revenue-and-profit-across-the-industry:elasticity}).

\begin{remark}[What the elasticity is and is not]\label{rem:in:revenue-and-profit-across-the-industry:elasticity}
Seven years are seven observations of the market; the standard errors assume independent residuals, which the shared
years do not give, so they are too small if anything. The elasticity describes the filers in the panel over one period
with one large volatility event, 2020. It is a description of how the kinds of firm differed, not a forecast: the next
volatile year may be driven by rates or credit, where a different set of firms earns.
\end{remark}

For a job seeker the regression has a plain reading. A market maker's income per person, and with it the variable part
of its pay (chapter~13), rises and falls with the market; an exchange's, a venue's or a bank's does much less.
The two asset managers, whose fees follow their funds' assets and results, show no clear link either.

\section{Who is missing: survivorship and disclosure bias}

\begin{definition}[Disclosure bias]\label{def:in:revenue-and-profit-across-the-industry:disclosure}
\emph{Disclosure bias}\index{disclosure bias} is the bias in a statistic over firms that arises because the firms that
publish the data differ systematically from those that do not: listed, regulated or registered where accounts are
public, rather than private and registered where they are not.
\end{definition}

The table's high end is set by the firms that publish, and the industry's high end by firms that do not. The largest
private market makers of chapter~2 file no consolidated income statement anywhere public: their US broker-dealers file
balance sheets, their UK subsidiaries file entity accounts that chapter~11 showed cannot be read as the firm's, and
press reports of their revenues are, by this book's rule, context and never data. A reader who wants the industry's
\omterm{def:in:revenue-and-profit-across-the-industry:perhead}{revenue per head} at the top has no filed number to use, and should say so rather than fill the cell.

Survivorship bias (Book~7, chapter~3) compounds it. The panel is built from firms that exist and file in 2025; a firm
that was acquired, delisted or closed during the period leaves no row, and a firm that closed because its revenue
collapsed takes its worst years with it. Both biases push in known directions: \omterm{def:in:revenue-and-profit-across-the-industry:disclosure}{disclosure bias} leaves out the most profitable private firms,
so the table understates the top; survivorship leaves out the failures, so the averages flatter the survivors. What the
table measures reliably is the shape of the industry's income, kind by kind, among firms that can be checked.

\section{Tutorial: the industry table}\label{tut:in:revenue-and-profit-across-the-industry:tut}

\begin{tutorial}
\textbf{Goal.} Build the firm-year panel from sources, normalise it, and estimate the elasticities.
\textbf{End state:} the dated table,
\cref{fig:in:revenue-and-profit-across-the-industry:ranges,fig:in:revenue-and-profit-across-the-industry:series,fig:in:revenue-and-profit-across-the-industry:elasticity}.
\begin{tutsteps}
\item \textbf{Sources.} \texttt{data/industry/headcounts.csv} holds each headcount with its filing and a quoted phrase;
\texttt{revenue\_manual.csv} the filers read from reports and the register. The chapter's script
\texttt{in\_panel\_derive.py}, run with a scratch directory, reads SEC company facts with \texttt{firm.filings} and
Book~16's tables, and writes \texttt{filings\_perhead.csv}; the tests read only the committed files.
\item \textbf{Normalise.} \texttt{firm.industrypnl.per\_head(rows, fx, cpi, 2025)} converts at the ECB annual rates of
chapter~7 and deflates with \texttt{data/industry/cpi\_usa.csv}.
\item \textbf{Fit.} \texttt{fe\_fit(ph, vix)} builds the design matrix of firm dummies and kind-specific slopes
(\cref{lst:in:revenue-and-profit-across-the-industry:fe}).
\omcode{code/firm/industrypnl/firm_industrypnl.py}{88}{113}{Firm fixed effects and one slope per kind of firm, by least
squares, with classical standard errors.}\label{lst:in:revenue-and-profit-across-the-industry:fe}
\item \textbf{Read.} \texttt{in\_industry.firm\_ranges}, \texttt{spread} and \texttt{mm\_ratios} give the numbers of
sections~2 and~3.
\end{tutsteps}
\end{tutorial}

\textbf{What to change next.} Replace the VIX with a measure closer to each kind's business (a bond volatility index
for the venues that trade bonds, crypto prices for the crypto exchange) and compare the fits; cluster the standard
errors by year.

\section{Build: the industry panel}\label{bld:in:revenue-and-profit-across-the-industry:industrypnl}

\begin{build}
\bfield{Purpose} One sourced firm-year panel of revenue and headcount, normalised, for this chapter and for chapter~14's
comparison of pay with revenue.
\bfield{Interface} \texttt{firm.industrypnl}: \texttt{Row}; \texttt{load\_panel}, \texttt{load\_fx},
\texttt{load\_cpi}; \texttt{real\_usd(row, fx, cpi, base)}; \texttt{per\_head(rows, fx, cpi, base)};
\texttt{ranges(ph)}; \texttt{ratio(ph, firm, y1, y0)}; \texttt{fe\_fit(ph, x, kinds)}. NumPy; wraps
\texttt{firm.bankmix}'s rates, and the panel is built with \texttt{firm.filings}.
\bfield{Rules} Rows without a headcount are dropped, never filled; revenue is in the filer's currency until converted;
the regression uses whole firms with point headcounts only; standard errors are classical and reported as indicative.
\bfield{Acceptance tests} \texttt{code/firm/industrypnl/tests/}: on a panel built exactly from known slopes the fit
recovers them with zero residual; with noise the estimates fall within four standard errors; conversion and
deflation match hand arithmetic; a row without a headcount is not loaded.
\bfield{Stretch} Standard errors clustered by year or by firm; a random-effects version; operating \omterm{def:in:revenue-and-profit-across-the-industry:perhead}{profit per head}
where the lines can be aligned.
\end{build}

\begin{omsources}
\item Forms 10-K and SEC company facts of CME Group, Intercontinental Exchange, Nasdaq, Cboe Global Markets,
MarketAxess, Tradeweb, Goldman Sachs, Morgan Stanley and Interactive Brokers, 2019--2025.
\item Flow Traders, Annual Reports 2019 and 2021 and results releases for 2023 and 2025.
\item Hudson River Trading Europe Ltd., accounts for 2023 and 2025 (Companies House).
\item OECD, consumer price indices; ECB, euro reference rates.
\item Book~16's derived tables for Virtu Financial, Man Group and the VIX; chapters~3, 10 and~11 for Optiver, Coinbase
and Quadrature Capital.
\end{omsources}

\section{Exercises}

\begin{exercise}[$\star$]\label{exo:in:revenue-and-profit-across-the-industry:1}
An exchange group reports total revenues of \$12\,640 million, operating expenses of \$5\,002 million and operating
income of \$4\,929 million. What are its revenues less transaction-based expenses, and how large were those expenses?
\end{exercise}

\begin{exercise}[$\star$]\label{exo:in:revenue-and-profit-across-the-industry:2}
From the dated table, which filers' \omterm{def:in:revenue-and-profit-across-the-industry:perhead}{revenue per head} was below \$1 million in their latest year, and what is the ratio
of the highest whole-firm figure to the lowest?
\end{exercise}

\begin{exercise}[$\star$]\label{exo:in:revenue-and-profit-across-the-industry:3}
Compute CME Group's operating income per head in 2025 and Optiver's net \omterm{def:in:revenue-and-profit-across-the-industry:perhead}{profit per head} in 2024 in euros.
\end{exercise}

\begin{exercise}[$\star\star$]\label{exo:in:revenue-and-profit-across-the-industry:4}
Nasdaq's \omterm{def:in:revenue-and-profit-across-the-industry:perhead}{revenue per head} fell to its lowest in 2023. Using its 10-K, give the explanation that is not about the market.
\end{exercise}

\begin{exercise}[$\star\star$]\label{exo:in:revenue-and-profit-across-the-industry:5}
With an elasticity of 1.19, by how much should a market maker's \omterm{def:in:revenue-and-profit-across-the-industry:perhead}{revenue per head} change from a year with a mean VIX of
16.85 to one of 29.25? Compare with Virtu's and Flow Traders' actual ratios.
\end{exercise}

\begin{exercise}[$\star\star$]\label{exo:in:revenue-and-profit-across-the-industry:6}
Why is Hudson River Trading's UK subsidiary kept out of the regression, although its figures are sourced?
\end{exercise}

\begin{exercise}[$\star\star\star$]\label{exo:in:revenue-and-profit-across-the-industry:7}
\emph{Coding.} Refit the regression without the crypto exchange and the broker. How do the market makers' and the
exchanges' elasticities and standard errors change, and why so little?
\end{exercise}

\begin{exercise}[$\star\star\star$]\label{exo:in:revenue-and-profit-across-the-industry:8}
\emph{Find the flaw.} ``Filed numbers show that the industry's \omterm{def:in:revenue-and-profit-across-the-industry:perhead}{revenue per head} peaks at about \$13 million, at a
systematic investment manager; market makers are well below that.''
\end{exercise}

\section{Problem: The Good Year}

\begin{problem}[{Weekend problem --- the good year}]\label{pb:in:revenue-and-profit-across-the-industry:1}
Two offers: one from a listed market maker, one from an exchange group. A friend says the market maker ``has good
years''. How good, and how often?

\medskip\noindent\textbf{Part I --- The table.}
\begin{enumerate}
\item Define \omterm{def:in:revenue-and-profit-across-the-industry:perhead}{revenue per head} and \omterm{def:in:revenue-and-profit-across-the-industry:perhead}{profit per head}, and list what must be reported with them.
\item Why use revenue after pass-through costs? Give the exchange group's example.
\item Name the bases on which filers state headcount and how each biases \omterm{def:in:revenue-and-profit-across-the-industry:perhead}{revenue per head}.
\item How are figures in euros and pounds, and from different years, made comparable?
\item How many firms and firm-years are in the panel, and where do they come from?
\end{enumerate}

\medskip\noindent\textbf{Part II --- What it hides.}
\begin{enumerate}[resume]
\item Give the lowest and highest whole-firm \omterm{def:in:revenue-and-profit-across-the-industry:perhead}{revenue per head}, with firm and year.
\item Give the ratio of highest to lowest year for Flow Traders, Virtu and CME Group.
\item Explain Nasdaq's and the Intercontinental Exchange's figures in 2023.
\item Why can a bank's markets division not be placed in the table?
\item What happened to the UK subsidiary's \omterm{def:in:revenue-and-profit-across-the-industry:perhead}{revenue per head} from 2022 to 2025, and why is it not the firm's?
\end{enumerate}

\medskip\noindent\textbf{Part III --- The regimes.}
\begin{enumerate}[resume]
\item Give the yearly mean VIX for 2019--2025.
\item Write the regression and say what the firm effects do.
\item Give each kind's elasticity and standard error.
\item Why is the crypto exchange's elasticity negative?
\item Give the market makers' 2020 to 2023 ratios of \omterm{def:in:revenue-and-profit-across-the-industry:perhead}{revenue per head}.
\end{enumerate}

\medskip\noindent\textbf{Part IV --- The verdict.}
\begin{enumerate}[resume]
\item State the \emph{named result}: the elasticity of \omterm{def:in:revenue-and-profit-across-the-industry:perhead}{revenue per head} to the mean VIX by kind of firm with standard
errors, and the range of the 2020 to 2023 ratio across the market makers.
\item Why are the standard errors optimistic?
\item Define \omterm{def:in:revenue-and-profit-across-the-industry:disclosure}{disclosure bias} and say which way it pushes the table's top.
\item How does survivorship bias push the averages?
\item In two sentences, answer the friend.
\end{enumerate}
\end{problem}

\section{Interview questions}

\begin{interviewq}[$\star$ \iqroles{trader}]\label{iq:in:revenue-and-profit-across-the-industry:1}
Why does a market maker earn more in a volatile year? Give two mechanisms.
\end{interviewq}

\begin{interviewq}[$\star$ \iqroles{bank, risk}]\label{iq:in:revenue-and-profit-across-the-industry:2}
An exchange reports total revenues and revenues less transaction-based expenses. Which do you use to compare it with a
bank, and why?
\end{interviewq}

\begin{interviewq}[$\star\star$ \iqroles{researcher}]\label{iq:in:revenue-and-profit-across-the-industry:3}
You regress log \omterm{def:in:revenue-and-profit-across-the-industry:perhead}{revenue per head} on log VIX with firm fixed effects over seven years. What does the fixed effect remove,
and what would a year fixed effect do to your slope?
\end{interviewq}

\begin{interviewq}[$\star\star$ \iqroles{researcher, mle}]\label{iq:in:revenue-and-profit-across-the-industry:4}
Your panel has fifteen firms and seven years. How would you compute honest standard errors?
\end{interviewq}

\begin{interviewq}[$\star\star$ \iqroles{developer, mle}]\label{iq:in:revenue-and-profit-across-the-industry:5}
Design the storage for a panel whose every cell must carry its source and whose sources restate. What are the keys?
\end{interviewq}

\begin{interviewq}[$\star\star\star$ \iqroles{researcher, trader}]\label{iq:in:revenue-and-profit-across-the-industry:6}
The largest market makers do not publish. How would you estimate their \omterm{def:in:revenue-and-profit-across-the-industry:perhead}{revenue per head}, and what would you refuse to
conclude?
\end{interviewq}
